\documentclass[11pt,a4paper]{article}
\usepackage[T1]{fontenc}
\usepackage{amsmath,amssymb}
\usepackage{geometry}
\usepackage{xcolor}
\usepackage{graphicx}
\usepackage{enumitem}
\usepackage{booktabs}
\usepackage{longtable}
\usepackage{array}
\usepackage[expansion=false]{microtype}
\usepackage{fancyhdr}
\usepackage{titlesec}
\usepackage{needspace}
\usepackage{etoolbox}
\usepackage{tocloft}
\usepackage{fvextra}

\usepackage{calc}

% Pandoc's Shaded + Highlighting environments (simple fallback)
\usepackage{framed}
\definecolor{shadecolor}{HTML}{f5f5f5}
\newenvironment{Shaded}{\begin{snugshade}\small}{\end{snugshade}}
% Stub style definitions used by Highlighting
\newcommand{\ImportTok}[1]{#1}
\newcommand{\KeywordTok}[1]{\textbf{#1}}
\newcommand{\DataTypeTok}[1]{#1}
\newcommand{\DecValTok}[1]{#1}
\newcommand{\BaseNTok}[1]{#1}
\newcommand{\FloatTok}[1]{#1}
\newcommand{\ConstantTok}[1]{#1}
\newcommand{\CharTok}[1]{#1}
\newcommand{\SpecialCharTok}[1]{#1}
\newcommand{\StringTok}[1]{\textit{#1}}
\newcommand{\VerbatimStringTok}[1]{\textit{#1}}
\newcommand{\SpecialStringTok}[1]{#1}
\newcommand{\CommentTok}[1]{\textit{\textcolor{gray}{#1}}}
\newcommand{\DocumentationTok}[1]{#1}
\newcommand{\AnnotationTok}[1]{#1}
\newcommand{\CommentVarTok}[1]{#1}
\newcommand{\OtherTok}[1]{#1}
\newcommand{\FunctionTok}[1]{#1}
\newcommand{\VariableTok}[1]{#1}
\newcommand{\ControlFlowTok}[1]{\textbf{#1}}
\newcommand{\OperatorTok}[1]{#1}
\newcommand{\BuiltInTok}[1]{#1}
\newcommand{\ExtensionTok}[1]{#1}
\newcommand{\PreprocessorTok}[1]{#1}
\newcommand{\AttributeTok}[1]{#1}
\newcommand{\RegionMarkerTok}[1]{#1}
\newcommand{\InformationTok}[1]{#1}
\newcommand{\WarningTok}[1]{#1}
\newcommand{\AlertTok}[1]{#1}
\newcommand{\ErrorTok}[1]{#1}
\newcommand{\NormalTok}[1]{#1}
\usepackage{fancyvrb}
\DefineVerbatimEnvironment{Highlighting}{Verbatim}{commandchars=\\\{\},frame=none}

\usepackage{newunicodechar}
\newunicodechar{Δ}{\ensuremath{\Delta}}
\newunicodechar{Φ}{\ensuremath{\Phi}}
\newunicodechar{ω}{\ensuremath{\omega}}
\newunicodechar{ε}{\ensuremath{\varepsilon}}
\newunicodechar{π}{\ensuremath{\pi}}
\newunicodechar{γ}{\ensuremath{\gamma}}
\newunicodechar{θ}{\ensuremath{\theta}}
\newunicodechar{τ}{\ensuremath{\tau}}
\newunicodechar{σ}{\ensuremath{\sigma}}
\newunicodechar{ρ}{\ensuremath{\rho}}
\newunicodechar{α}{\ensuremath{\alpha}}
\newunicodechar{β}{\ensuremath{\beta}}
\newunicodechar{ν}{\ensuremath{\nu}}
\newunicodechar{μ}{\ensuremath{\mu}}
\newunicodechar{→}{\ensuremath{\rightarrow}}
\newunicodechar{←}{\ensuremath{\leftarrow}}
\newunicodechar{↔}{\ensuremath{\leftrightarrow}}
\newunicodechar{≥}{\ensuremath{\geq}}
\newunicodechar{≤}{\ensuremath{\leq}}
\newunicodechar{≠}{\ensuremath{\neq}}
\newunicodechar{⁻}{\ensuremath{{}^-}}
\newunicodechar{¹}{\ensuremath{{}^1}}
\newunicodechar{²}{\ensuremath{{}^2}}
\newunicodechar{³}{\ensuremath{{}^3}}
\newunicodechar{⁴}{\ensuremath{{}^4}}
\newunicodechar{⁶}{\ensuremath{{}^6}}
\newunicodechar{⁸}{\ensuremath{{}^8}}
\newunicodechar{·}{\ensuremath{\cdot}}
\newunicodechar{×}{\ensuremath{\times}}
\newunicodechar{∈}{\ensuremath{\in}}
\newunicodechar{ℚ}{\ensuremath{\mathbb{Q}}}
\newunicodechar{ℝ}{\ensuremath{\mathbb{R}}}
\newunicodechar{⅓}{\ensuremath{\tfrac{1}{3}}}
\newunicodechar{└}{\textbar}
\newunicodechar{─}{-}
\newunicodechar{┐}{\textbar}
\newunicodechar{┌}{\textbar}
\newunicodechar{┘}{\textbar}
\newunicodechar{│}{\textbar}
\newunicodechar{├}{\textbar}
\newunicodechar{┤}{\textbar}
\newunicodechar{┬}{-}
\newunicodechar{┴}{-}
\newunicodechar{┼}{+}
\newunicodechar{▼}{v}
\newunicodechar{▲}{\^{}}
\newunicodechar{▶}{>}
\newunicodechar{↑}{\ensuremath{\uparrow}}
\newunicodechar{↓}{\ensuremath{\downarrow}}
\newunicodechar{⟶}{\ensuremath{\longrightarrow}}
\newunicodechar{⟵}{\ensuremath{\longleftarrow}}

\newunicodechar{₁}{\ensuremath{{}_1}}
\newunicodechar{₂}{\ensuremath{{}_2}}
\newunicodechar{₃}{\ensuremath{{}_3}}
\newunicodechar{₄}{\ensuremath{{}_4}}
\newunicodechar{₆}{\ensuremath{{}_6}}
\newunicodechar{⁰}{\ensuremath{{}^0}}
\newunicodechar{⁵}{\ensuremath{{}^5}}
\newunicodechar{⁷}{\ensuremath{{}^7}}
\newunicodechar{⁹}{\ensuremath{{}^9}}
\newunicodechar{π}{\ensuremath{\pi}}
\newunicodechar{ε}{\ensuremath{\varepsilon}}

\usepackage{xltabular}
\usepackage{tabularx}

\definecolor{acsprimary}{HTML}{1a237e}
\definecolor{acssecondary}{HTML}{3949ab}
\definecolor{acsaccent}{HTML}{c62828}
\definecolor{acsgold}{HTML}{b8860b}
\definecolor{acsgreen}{HTML}{2e7d32}
\definecolor{acsslate}{HTML}{37474f}
\definecolor{acslight}{HTML}{e8eaf6}
\definecolor{acsrule}{HTML}{9fa8da}
\definecolor{codeback}{HTML}{f5f5f5}

\usepackage[labelfont={bf,small,color=acsprimary!90!black},textfont={small,it}]{caption}
\usepackage{hyperref}
\usepackage{bookmark}

\geometry{top=1.15in,bottom=1.15in,left=1.2in,right=1.2in,
          headheight=16pt,headsep=0.25in,footskip=0.45in}

\hypersetup{colorlinks=true,linkcolor=acsprimary,citecolor=acssecondary,
            urlcolor=acsaccent,pdfpagemode=UseOutlines,bookmarksopen=true}

\pagestyle{fancy}
\fancyhf{}
\renewcommand{\headrulewidth}{0.4pt}
\renewcommand{\footrulewidth}{0pt}
\fancyhead[L]{\small\color{acsslate}\itshape\leftmark}
\fancyhead[R]{\small\color{acsslate}\textit{ACS Deterministic AI Stack PDR}}
\fancyfoot[C]{\small\color{acsslate}\thepage}
\renewcommand{\sectionmark}[1]{\markboth{\S\thesection.\ #1}{}}

\titleformat{\section}
  {\needspace{6\baselineskip}\normalfont\Large\bfseries\color{acsprimary}}
  {\thesection}{0.7em}{}
  [\vspace{-2pt}{\color{acsrule}\rule{\textwidth}{0.4pt}}]
\titleformat{\subsection}
  {\needspace{5\baselineskip}\normalfont\large\bfseries\color{acssecondary}}
  {\thesubsection}{0.7em}{}
\titleformat{\subsubsection}
  {\needspace{4\baselineskip}\normalfont\normalsize\bfseries\color{acsslate}}
  {\thesubsubsection}{0.7em}{}
\titlespacing*{\section}{0pt}{2.4ex plus 0.9ex minus 0.4ex}{1.4ex plus 0.3ex}
\titlespacing*{\subsection}{0pt}{2.0ex plus 0.7ex minus 0.3ex}{0.9ex plus 0.2ex}
\titlespacing*{\subsubsection}{0pt}{1.6ex plus 0.5ex minus 0.2ex}{0.7ex}

\setlist[itemize]{leftmargin=1.3em,itemsep=2pt,topsep=3pt,parsep=0pt}
\setlist[enumerate]{leftmargin=1.5em,itemsep=2pt,topsep=3pt,parsep=0pt}

\setcounter{tocdepth}{2}
\setlength{\cftbeforesecskip}{2pt}
\setlength{\cftbeforesubsecskip}{0.5pt}
\providecommand{\tightlist}{\setlength{\itemsep}{0pt}\setlength{\parskip}{0pt}}
\widowpenalty=10000
\clubpenalty=10000
\displaywidowpenalty=10000
\begin{document}
\begin{titlepage}
\thispagestyle{empty}
\noindent{\color{acsrule}\rule{\textwidth}{0.4pt}}
\vspace{6pt}
\noindent{\footnotesize\color{acsslate}\hfill\textit{ACS Deterministic AI Stack PDR} \quad|\quad \textit{Architectural PDR} \quad|\quad April 2026}
\vspace{6pt}
\noindent{\color{acsrule}\rule{\textwidth}{0.4pt}}
\vfill
\begin{center}
  {\Huge\bfseries\color{acsprimary} ACS Deterministic AI Stack}\\[0.4em]
  {\Large\color{acsslate} Preliminary Design Review \& Technical Blueprint}\\[1.2em]
  {\normalsize\itshape\color{acsslate}
    A bracket-governed, byte-reproducible architecture\\
    for modelling the Asymmetric Codependent Systems framework}\\[2.0em]
  {\large\bfseries Bradley Wallace}\\[0.25em]
  {\small\itshape\color{acsslate} Independent Researcher}
\end{center}
\vspace{0.4in}
\begin{center}
\begin{minipage}{0.88\textwidth}
\begin{center}
  {\bfseries\color{acsprimary}\large Executive Summary}\\[0.3em]
  {\color{acsrule}\rule{1.2in}{0.3pt}}
\end{center}
\vspace{0.5em}
\small\noindent
This document specifies a deterministic AI stack governed by the
Asymmetric Codependent Systems (ACS) framework. Every generative step
is a Lie bracket between two typed fields (Form and Function); outputs
are validated by a closure criterion; convergence is detected by
information balance ($\Delta\mathcal{I} \to 0$); compositional depth
is capped at the third BCH order via the Jacobi identity. The stack
is specified in five layers, four deployable services, and a phased
roadmap from foundations through extension to the physics frontier.
Phase~1 targets full deployment of the 76-script ACS verification suite
under the new governance within a six-month horizon.
\end{minipage}
\end{center}
\vfill
\noindent{\color{acsrule}\rule{\textwidth}{0.4pt}}
\end{titlepage}
\clearpage
\pagenumbering{roman}
\setcounter{page}{1}
\markboth{Contents}{}
\tableofcontents
\clearpage
\pagenumbering{arabic}
\setcounter{page}{1}

\hypertarget{governing-principles-acs-ai-mapping}{%
\section{Governing Principles --- ACS → AI
Mapping}\label{governing-principles-acs-ai-mapping}}

The physics framework maps onto software architecture through a set of
structural correspondences. These are not analogies; they are direct
mappings because the ACS governs information flow in \emph{any}
codependent system, whether physical or computational.

\begin{longtable}[]{@{}
  >{\raggedright\arraybackslash}p{(\columnwidth - 2\tabcolsep) * \real{0.5652}}
  >{\raggedright\arraybackslash}p{(\columnwidth - 2\tabcolsep) * \real{0.4348}}@{}}
\toprule\noalign{}
\begin{minipage}[b]{\linewidth}\raggedright
ACS Physics
\end{minipage} & \begin{minipage}[b]{\linewidth}\raggedright
AI Stack
\end{minipage} \\
\midrule\noalign{}
\endhead
\bottomrule\noalign{}
\endlastfoot
Form field \texttt{F} (vierbein \texttt{e}) & Typed state schema --- the
``what'' \\
Function field \texttt{Φ} (connection \texttt{ω}) & Operator algebra ---
the ``how'' \\
Lie bracket \texttt{{[}F,\ Φ{]}} & Composition core --- generative
step \\
BCH order 1 (direct coupling) & I/O layer --- input processing \\
BCH order 2 (curvature \texttt{{[}f,g{]}}) & Operator composition ---
transformations \\
BCH order 3 (holonomy \texttt{{[}{[}f,g{]},·{]}}) & Emergent closure ---
final output \\
Jacobi identity (truncates at 3) & Bounded-depth guarantee \\
Closure attractor (\texttt{D\ \textless{}\ 10⁻¹⁴}) & Convergence
validator \\
Information balance \texttt{ΔI\ =\ 0} & Halting criterion \\
Constraint-attractor cycle & Governance feedback loop \\
Torsion hierarchy \texttt{0:1:4} & Operator priority tiers \\
Three generations & Three irreducible output channels \\
7-parameter boundary & Explicit configuration surface \\
\end{longtable}

\textbf{The invariant that makes this deterministic.} For every
generative step, there exists a closed-form algebraic expression in the
operator algebra that produces the output exactly. There is no sampling.
There is no temperature. The ``creativity'' comes from the bracket
itself, whose structure is determined once the two fields are fixed.

\begin{center}\rule{0.5\linewidth}{0.5pt}\end{center}

\hypertarget{system-architecture}{%
\section{System Architecture}\label{system-architecture}}

\hypertarget{high-level-topology}{%
\subsection{High-Level Topology}\label{high-level-topology}}

\begin{verbatim}
+-------------------------------------------------------------------+
|                    CLIENT / USER INTERFACE                         |
|           (CLI · REST API · Jupyter · Verification Report)         |
+-----------------------------+-------------------------------------+
                              |
+-----------------------------v-------------------------------------+
|                    GOVERNANCE LAYER                                |
|   Policy Engine · Closure Validator · ΔI Monitor · Audit Log      |
+-----------------------------+-------------------------------------+
                              |
+-----------------------------v-------------------------------------+
|                   COMPOSITION CORE (Bracket Engine)                |
|    BCH-3 Truncated Executor · Jacobi Enforcer · Holonomy Tracker  |
+--------------+--------------------------+-------------------------+
               |                          |
   +-----------v----------+   +----------v-----------+
   |    FORM LAYER        |   |    FUNCTION LAYER    |
   |  Typed state schema  |   |   Operator algebra   |
   |  (Pydantic + SymPy)  |   |  (registered brackets)|
   +-----------+----------+   +----------+-----------+
               |                          |
+--------------v--------------------------v-------------------------+
|                    KERNEL — Symbolic + Numerical                  |
|   SymPy (exact Q) · mpmath (arbitrary precision) · NumPy (fp64)   |
+-----------------------------+-------------------------------------+
                              |
+-----------------------------v-------------------------------------+
|                    PERSISTENCE & PROVENANCE                        |
|   Content-addressed store (SHA-256) · DAG journal · Immutable log |
+-------------------------------------------------------------------+
\end{verbatim}

\hypertarget{the-five-architectural-layers}{%
\subsection{The Five Architectural
Layers}\label{the-five-architectural-layers}}

\textbf{Layer 0: Kernel (numerical foundation).} Provides three
interchangeable arithmetic backends --- SymPy for exact rational
arithmetic (the only mode admissible for theorem-grade claims), mpmath
for arbitrary-precision floating-point, and NumPy for standard double
precision. All three must produce bit-identical results for any
computation representable in their common domain; this is the
\emph{kernel invariant}, checked at every build.

\textbf{Layer 1: Form and Function (typed data plane).} All data in the
stack is classified as either Form (state, configuration, schemas) or
Function (operators, transformations, brackets). This classification is
enforced at the type system level using Pydantic for runtime validation
and mypy for static checking. A value cannot simultaneously inhabit both
categories; this is the \emph{codependence invariant}.

\textbf{Layer 2: Composition Core (the Bracket Engine).} The generative
engine. Takes a Form \texttt{F} and a Function \texttt{Φ} and produces
their bracket \texttt{{[}F,\ Φ{]}} truncated at BCH order 3. Every
composition is recorded in the DAG journal with cryptographic
provenance. The engine is stateless between brackets --- all state lives
in Layer 1.

\textbf{Layer 3: Governance (the policy plane).} Enforces the closure
criterion, information-balance halting, and the constraint-attractor
cycle. Every output from Layer 2 passes through Layer 3 before reaching
the client. Rejected outputs are returned with a machine-readable
violation report.

\textbf{Layer 4: Interface (client plane).} CLI, REST API, Jupyter
kernel, and verification reports. The API contract is minimal and
explicit; there is no hidden state, no conversational memory, no
implicit context. Every call is idempotent.

\hypertarget{dependency-direction-and-invariants}{%
\subsection{Dependency Direction and
Invariants}\label{dependency-direction-and-invariants}}

\begin{itemize}
\tightlist
\item
  Dependencies flow downward only. Layer 3 depends on Layer 2, never the
  reverse. This prevents governance bypass.
\item
  Layers 0 and 1 are pure (no side effects). Layer 2 is pure modulo the
  DAG journal append. Only Layers 3 and 4 may have external side
  effects.
\item
  All inter-layer communication uses immutable, hashed message objects.
\end{itemize}

\begin{center}\rule{0.5\linewidth}{0.5pt}\end{center}

\hypertarget{component-specifications}{%
\section{Component Specifications}\label{component-specifications}}

\hypertarget{the-form-layer}{%
\subsection{The Form Layer}\label{the-form-layer}}

\textbf{Purpose.} Represent all typed state in the system. Corresponds
to the vierbein \texttt{e} in the physics: the field that encodes
``shape''.

\textbf{Interface.}

\begin{Shaded}
\begin{Highlighting}[]
\ImportTok{from}\NormalTok{ pydantic }\ImportTok{import}\NormalTok{ BaseModel}
\ImportTok{from}\NormalTok{ typing }\ImportTok{import}\NormalTok{ TypeVar, Generic}
\ImportTok{from}\NormalTok{ hashlib }\ImportTok{import}\NormalTok{ sha256}

\NormalTok{T }\OperatorTok{=}\NormalTok{ TypeVar(}\StringTok{\textquotesingle{}T\textquotesingle{}}\NormalTok{)}

\KeywordTok{class}\NormalTok{ Form(BaseModel, Generic[T]):}
    \CommentTok{"""A Form field carrying structural information."""}
\NormalTok{    name: }\BuiltInTok{str}
\NormalTok{    content: T}
\NormalTok{    schema\_version: }\BuiltInTok{str}
\NormalTok{    provenance: }\BuiltInTok{str}  \CommentTok{\# SHA{-}256 of the generating bracket chain, or "initial"}

    \KeywordTok{class}\NormalTok{ Config:}
\NormalTok{        frozen }\OperatorTok{=} \VariableTok{True}  \CommentTok{\# Forms are immutable}

    \AttributeTok{@property}
    \KeywordTok{def}\NormalTok{ content\_hash(}\VariableTok{self}\NormalTok{) }\OperatorTok{{-}\textgreater{}} \BuiltInTok{str}\NormalTok{:}
        \CommentTok{"""Content{-}addressed identity."""}
\NormalTok{        canonical }\OperatorTok{=} \VariableTok{self}\NormalTok{.json(sort\_keys}\OperatorTok{=}\VariableTok{True}\NormalTok{).encode()}
        \ControlFlowTok{return}\NormalTok{ sha256(canonical).hexdigest()}

    \KeywordTok{def} \FunctionTok{\_\_or\_\_}\NormalTok{(}\VariableTok{self}\NormalTok{, other: }\StringTok{\textquotesingle{}Function\textquotesingle{}}\NormalTok{) }\OperatorTok{{-}\textgreater{}} \StringTok{\textquotesingle{}Form\textquotesingle{}}\NormalTok{:}
        \CommentTok{"""Syntactic sugar: F | Φ computes [F, Φ] via the Bracket Engine."""}
        \ImportTok{from}\NormalTok{ acs.engine }\ImportTok{import}\NormalTok{ bracket}
        \ControlFlowTok{return}\NormalTok{ bracket(}\VariableTok{self}\NormalTok{, other)}
\end{Highlighting}
\end{Shaded}

\textbf{Invariants.}

\begin{itemize}
\tightlist
\item
  \texttt{I-F1} (Immutability): A Form, once constructed, is never
  mutated.
\item
  \texttt{I-F2} (Content-addressed): Two Forms with equal
  \texttt{content} and \texttt{schema\_version} produce identical
  \texttt{content\_hash}.
\item
  \texttt{I-F3} (Type safety): \texttt{content} is validated against
  \texttt{schema\_version} at construction time; construction fails on
  schema mismatch.
\end{itemize}

\textbf{Standard Form types shipped with the kernel.}

\begin{itemize}
\tightlist
\item
  \texttt{AlgebraElement} --- element of a Lie algebra (e.g., a matrix
  in sl(4,ℝ))
\item
  \texttt{TensorState} --- typed multilinear tensor
\item
  \texttt{Spectrum} --- a finite spectrum (eigenvalues + multiplicities)
\item
  \texttt{Measurement} --- a physical observable with uncertainty
\item
  \texttt{Configuration} --- a frozen parameter bundle
\end{itemize}

\hypertarget{the-function-layer}{%
\subsection{The Function Layer}\label{the-function-layer}}

\textbf{Purpose.} Represent all operators/transformations. Corresponds
to the spin connection \texttt{ω}: the field that generates change.

\textbf{Interface.}

\begin{Shaded}
\begin{Highlighting}[]
\ImportTok{from}\NormalTok{ abc }\ImportTok{import}\NormalTok{ ABC, abstractmethod}

\KeywordTok{class}\NormalTok{ Function(ABC):}
    \CommentTok{"""A Function is an operator acting on Forms."""}
\NormalTok{    name: }\BuiltInTok{str}
\NormalTok{    order: }\BuiltInTok{int}  \CommentTok{\# BCH order at which this operator acts (1, 2, or 3)}

    \AttributeTok{@abstractmethod}
    \KeywordTok{def} \BuiltInTok{apply}\NormalTok{(}\VariableTok{self}\NormalTok{, form: Form) }\OperatorTok{{-}\textgreater{}}\NormalTok{ Form:}
        \CommentTok{"""The base action. Must be deterministic and side{-}effect{-}free."""}

    \AttributeTok{@abstractmethod}
    \KeywordTok{def}\NormalTok{ bracket\_with(}\VariableTok{self}\NormalTok{, other: }\StringTok{\textquotesingle{}Function\textquotesingle{}}\NormalTok{) }\OperatorTok{{-}\textgreater{}} \StringTok{\textquotesingle{}Function\textquotesingle{}}\NormalTok{:}
        \CommentTok{"""The [self, other] composition. Returns a new Function."""}

    \AttributeTok{@abstractmethod}
    \KeywordTok{def}\NormalTok{ killing\_form(}\VariableTok{self}\NormalTok{, other: }\StringTok{\textquotesingle{}Function\textquotesingle{}}\NormalTok{) }\OperatorTok{{-}\textgreater{}}\NormalTok{ Rational:}
        \CommentTok{"""K(self, other) — the Killing form value. Exact rational."""}
\end{Highlighting}
\end{Shaded}

\textbf{Invariants.}

\begin{itemize}
\tightlist
\item
  \texttt{I-Φ1} (Determinism): \texttt{Φ.apply(F)} depends only on
  \texttt{Φ} and \texttt{F}.
\item
  \texttt{I-Φ2} (Bracket closure): \texttt{Φ₁.bracket\_with(Φ₂)} must
  live in the same algebra as \texttt{Φ₁} and \texttt{Φ₂}.
\item
  \texttt{I-Φ3} (Jacobi): For any three Functions,
  \texttt{{[}{[}Φ₁,Φ₂{]},Φ₃{]}\ +\ {[}{[}Φ₂,Φ₃{]},Φ₁{]}\ +\ {[}{[}Φ₃,Φ₁{]},Φ₂{]}\ =\ 0}
  (verified at registration time).
\item
  \texttt{I-Φ4} (Order tagging): Every Function declares its BCH order.
  Attempting to compose Functions such that the result would exceed
  order 3 raises \texttt{JacobiTruncationError}.
\end{itemize}

\textbf{Standard Function types.}

\begin{itemize}
\tightlist
\item
  \texttt{BracketOperator} --- a Lie algebra generator acting by
  commutator
\item
  \texttt{ProjectionOperator} --- orthogonal projection onto a subspace
\item
  \texttt{KillingOperator} --- the Killing form K(X, Y)
\item
  \texttt{DifferentialOperator} --- covariant derivative
\item
  \texttt{ClosureTest} --- computes the closure defect \texttt{D(V)}
\end{itemize}

\hypertarget{the-bracket-engine}{%
\subsection{The Bracket Engine}\label{the-bracket-engine}}

\textbf{Purpose.} The single generative primitive. Takes \texttt{F} and
\texttt{Φ} and produces \texttt{{[}F,\ Φ{]}} truncated at the third BCH
order.

\textbf{Interface.}

\begin{Shaded}
\begin{Highlighting}[]
\KeywordTok{def}\NormalTok{ bracket(}
\NormalTok{    f: Form,}
\NormalTok{    phi: Function,}
\NormalTok{    order: }\BuiltInTok{int} \OperatorTok{=} \DecValTok{3}\NormalTok{,}
\NormalTok{    validate\_closure: }\BuiltInTok{bool} \OperatorTok{=} \VariableTok{True}\NormalTok{,}
\NormalTok{) }\OperatorTok{{-}\textgreater{}}\NormalTok{ Form:}
    \CommentTok{"""}
\CommentTok{    Compute [F, Φ] truncated at the given BCH order.}

\CommentTok{    Parameters}
\CommentTok{    {-}{-}{-}{-}{-}{-}{-}{-}{-}{-}}
\CommentTok{    f : Form}
\CommentTok{        The Form field (state).}
\CommentTok{    phi : Function}
\CommentTok{        The Function field (operator).}
\CommentTok{    order : int in \{1, 2, 3\}}
\CommentTok{        BCH truncation order. Default 3.}
\CommentTok{        Order 1: direct coupling (Φ.apply(F)).}
\CommentTok{        Order 2: Lie bracket (curvature).}
\CommentTok{        Order 3: holonomy — double bracket.}
\CommentTok{    validate\_closure : bool}
\CommentTok{        If True, run the Closure Validator on the result.}

\CommentTok{    Returns}
\CommentTok{    {-}{-}{-}{-}{-}{-}{-}}
\CommentTok{    Form}
\CommentTok{        The output, with provenance pointing to (f, phi, order).}

\CommentTok{    Raises}
\CommentTok{    {-}{-}{-}{-}{-}{-}}
\CommentTok{    JacobiTruncationError : if order \textgreater{} 3.}
\CommentTok{    ClosureViolation : if validate\_closure and the result fails closure.}
\CommentTok{    BracketTypeError : if f and phi are not composable.}
\CommentTok{    """}
\end{Highlighting}
\end{Shaded}

\textbf{Execution semantics.}

\begin{enumerate}
\def\labelenumi{\arabic{enumi}.}
\tightlist
\item
  \textbf{Type check.} Verify that \texttt{phi} is applicable to
  \texttt{f}. This checks schema compatibility and the codependence
  invariant.
\item
  \textbf{Order-1 step.} Compute \texttt{r₁\ =\ Φ(F)}.
\item
  \textbf{Order-2 step.} If order ≥ 2, compute
  \texttt{r₂\ =\ {[}F,\ r₁{]}} using the algebra's bracket operation.
\item
  \textbf{Order-3 step.} If order = 3, compute
  \texttt{r₃\ =\ {[}{[}F,\ r₁{]},\ F{]}} and
  \texttt{r₃\textquotesingle{}\ =\ {[}{[}F,\ r₁{]},\ r₁{]}}, then sum.
\item
  \textbf{Journal append.} Write the triple
  \texttt{(content\_hash(f),\ phi.name,\ order)} and the output's
  \texttt{content\_hash} to the immutable DAG journal.
\item
  \textbf{Closure validation.} If enabled, run the output through the
  Closure Validator (§3.4). On failure, roll back the journal append
  (using the 2-phase-commit protocol specified in §5.3) and raise.
\item
  \textbf{Return.} Wrapped Form with provenance.
\end{enumerate}

\textbf{Invariants.}

\begin{itemize}
\tightlist
\item
  \texttt{I-B1} (Reproducibility): \texttt{bracket(f,\ phi,\ n)} called
  twice with byte-identical \texttt{f} and \texttt{phi} returns
  byte-identical Forms.
\item
  \texttt{I-B2} (BCH truncation): The engine never computes orders above
  3.
\item
  \texttt{I-B3} (Journal-before-return): No output is returned to the
  caller before its provenance is committed to the journal.
\item
  \texttt{I-B4} (Rollback): If any step fails, no partial state is
  visible outside the engine.
\end{itemize}

\hypertarget{the-closure-validator}{%
\subsection{The Closure Validator}\label{the-closure-validator}}

\textbf{Purpose.} Enforce the closure attractor principle --- outputs
must satisfy \texttt{D(output)\ \textless{}\ ε} where \texttt{D} is the
closure defect functional.

\textbf{Interface.}

\begin{Shaded}
\begin{Highlighting}[]
\ImportTok{from}\NormalTok{ decimal }\ImportTok{import}\NormalTok{ Decimal}

\AttributeTok{@dataclass}\NormalTok{(frozen}\OperatorTok{=}\VariableTok{True}\NormalTok{)}
\KeywordTok{class}\NormalTok{ ClosureResult:}
\NormalTok{    passed: }\BuiltInTok{bool}
\NormalTok{    defect: Decimal}
\NormalTok{    threshold: Decimal}
\NormalTok{    witness: Optional[}\BuiltInTok{str}\NormalTok{]  }\CommentTok{\# explanation of failure, if any}

\KeywordTok{def}\NormalTok{ validate\_closure(}
\NormalTok{    form: Form,}
\NormalTok{    subspace: AlgebraSubspace,}
\NormalTok{    threshold: Decimal }\OperatorTok{=}\NormalTok{ Decimal(}\StringTok{"1e{-}14"}\NormalTok{),}
\NormalTok{) }\OperatorTok{{-}\textgreater{}}\NormalTok{ ClosureResult:}
    \CommentTok{"""}
\CommentTok{    Compute D(subspace) for the algebra element in \textasciigrave{}form\textasciigrave{} and check}
\CommentTok{    whether it closes in \textasciigrave{}subspace\textasciigrave{} to within \textasciigrave{}threshold\textasciigrave{}.}

\CommentTok{    D(V) = sum\_\{i\textless{}j\} ||[T\_i, T\_j] {-} π\_V([T\_i, T\_j])||}
\CommentTok{         / sum\_\{i\textless{}j\} ||[T\_i, T\_j]||}

\CommentTok{    where \{T\_i\} is a basis of V and π\_V is orthogonal projection onto V.}
\CommentTok{    """}
\end{Highlighting}
\end{Shaded}

\textbf{Standard thresholds.}

\begin{itemize}
\tightlist
\item
  \texttt{Decimal("1e-14")} for exact rational computations (SymPy
  backend)
\item
  \texttt{Decimal("1e-10")} for arbitrary-precision (mpmath)
\item
  \texttt{Decimal("1e-6")} for double-precision (NumPy) --- advisory
  only, not admissible for theorem-grade claims
\end{itemize}

\textbf{Invariants.}

\begin{itemize}
\tightlist
\item
  \texttt{I-C1} (Monotonicity): If \texttt{D(V)\ \textless{}\ ε} with
  threshold \texttt{ε}, then for any
  \texttt{ε\textquotesingle{}\ \textgreater{}\ ε},
  \texttt{D(V)\ \textless{}\ ε\textquotesingle{}} also passes.
\item
  \texttt{I-C2} (Provenance): Every ClosureResult carries a witness
  string sufficient to reproduce the computation.
\end{itemize}

\hypertarget{the-ux3b4i-monitor}{%
\subsection{The ΔI Monitor}\label{the-ux3b4i-monitor}}

\textbf{Purpose.} Detect convergence via information balance. A
computation that reaches \texttt{ΔI\ →\ 0} has found its attractor and
should halt.

\textbf{Interface.}

\begin{Shaded}
\begin{Highlighting}[]
\KeywordTok{def}\NormalTok{ delta\_I(}
\NormalTok{    history: }\BuiltInTok{list}\NormalTok{[Form],}
\NormalTok{    window: }\BuiltInTok{int} \OperatorTok{=} \DecValTok{8}\NormalTok{,}
\NormalTok{) }\OperatorTok{{-}\textgreater{}}\NormalTok{ Decimal:}
    \CommentTok{"""}
\CommentTok{    Compute the net information asymmetry over the last \textasciigrave{}window\textasciigrave{} Forms}
\CommentTok{    in the computation history.}

\CommentTok{    ΔI = TE(F → Φ) {-} TE(Φ → F)}

\CommentTok{    where TE is transfer entropy estimated via the k{-}nearest{-}neighbour}
\CommentTok{    estimator (Kraskov{-}Stögbauer{-}Grassberger).}
\CommentTok{    """}

\KeywordTok{def}\NormalTok{ is\_converged(history: }\BuiltInTok{list}\NormalTok{[Form], epsilon: Decimal }\OperatorTok{=}\NormalTok{ Decimal(}\StringTok{"1e{-}6"}\NormalTok{)) }\OperatorTok{{-}\textgreater{}} \BuiltInTok{bool}\NormalTok{:}
    \CommentTok{"""True if |ΔI| \textless{} ε over the trailing window."""}
\end{Highlighting}
\end{Shaded}

\textbf{Halting policy.} The engine halts when any of the following
hold:

\begin{enumerate}
\def\labelenumi{\arabic{enumi}.}
\tightlist
\item
  \texttt{\textbar{}ΔI\textbar{}\ \textless{}\ ε} over the last 8 steps
  --- information-balance attractor
\item
  BCH order 3 has been reached and produced no new independent content
  (Jacobi truncation)
\item
  The closure defect has been below threshold for 3 consecutive steps
\item
  An explicit halt request is issued by the governance layer
\end{enumerate}

\hypertarget{the-governance-policy-engine}{%
\subsection{The Governance Policy
Engine}\label{the-governance-policy-engine}}

\textbf{Purpose.} Enforce the constraint-attractor cycle. Every output
must satisfy the current constraint set; outputs that satisfy become the
next constraint.

\textbf{Interface.}

\begin{Shaded}
\begin{Highlighting}[]
\AttributeTok{@dataclass}\NormalTok{(frozen}\OperatorTok{=}\VariableTok{True}\NormalTok{)}
\KeywordTok{class}\NormalTok{ Policy:}
\NormalTok{    name: }\BuiltInTok{str}
\NormalTok{    constraints: }\BuiltInTok{list}\NormalTok{[Constraint]}
\NormalTok{    attractor\_criteria: }\BuiltInTok{list}\NormalTok{[AttractorTest]}
\NormalTok{    version: }\BuiltInTok{str}

\KeywordTok{def}\NormalTok{ enforce\_policy(}
\NormalTok{    form: Form,}
\NormalTok{    policy: Policy,}
\NormalTok{) }\OperatorTok{{-}\textgreater{}}\NormalTok{ PolicyResult:}
    \CommentTok{"""}
\CommentTok{    Run \textasciigrave{}form\textasciigrave{} against all policy constraints and attractor tests.}
\CommentTok{    Returns either Approved(form) or Rejected(reasons, witness).}
\CommentTok{    """}

\KeywordTok{def}\NormalTok{ promote\_to\_constraint(}
\NormalTok{    form: Form,}
\NormalTok{    policy: Policy,}
\NormalTok{) }\OperatorTok{{-}\textgreater{}}\NormalTok{ Policy:}
    \CommentTok{"""}
\CommentTok{    If \textasciigrave{}form\textasciigrave{} has satisfied the attractor criteria, promote it to be}
\CommentTok{    a constraint for subsequent computations. Returns a new Policy}
\CommentTok{    with the form added to the constraint set.}

\CommentTok{    Implements the inversion arc: solution → next constraint.}
\CommentTok{    """}
\end{Highlighting}
\end{Shaded}

\textbf{Standard policies.}

\begin{itemize}
\tightlist
\item
  \texttt{TheoremGrade} --- requires SymPy backend, closure \textless{}
  1e-14, full Jacobi verification, no stochastic steps
\item
  \texttt{VerificationGrade} --- permits mpmath, closure \textless{}
  1e-10, Jacobi sampling
\item
  \texttt{ExploratoryGrade} --- permits NumPy, closure \textless{} 1e-6,
  advisory only
\end{itemize}

Every output is tagged with the grade of policy it satisfied.

\begin{center}\rule{0.5\linewidth}{0.5pt}\end{center}

\hypertarget{data-flow-and-execution-model}{%
\section{Data Flow and Execution
Model}\label{data-flow-and-execution-model}}

\hypertarget{the-canonical-computation-dag}{%
\subsection{The Canonical Computation
DAG}\label{the-canonical-computation-dag}}

Every computation in the stack is a directed acyclic graph whose:

\begin{itemize}
\tightlist
\item
  Nodes are Forms (typed immutable states)
\item
  Edges are Functions (operators labelled with BCH order)
\item
  Root nodes are Inputs (user-supplied or configuration)
\item
  Leaf nodes are Outputs (policy-approved Forms)
\end{itemize}

The DAG is explicit, serialisable, and content-addressed. Every edge is
labelled with the SHA-256 of
\texttt{(source\_hash,\ function\_name,\ order)}, making the entire
computation history a Merkle DAG.

\hypertarget{execution-phases}{%
\subsection{Execution Phases}\label{execution-phases}}

\begin{verbatim}
PHASE 1: INGEST
  input bytes → Pydantic validation → Form(provenance="initial")

PHASE 2: PLAN
  target specification → DAG of Form nodes and Function edges
  (this is a pure symbolic phase — no computation yet)

PHASE 3: EXECUTE
  topological walk of the DAG, dispatching each edge to the Bracket
  Engine. Cache hits (content_hash already in store) skip re-execution.

PHASE 4: VALIDATE
  every Form passes through Closure Validator and Policy Engine.
  Failures halt execution immediately with full provenance trace.

PHASE 5: CONVERGE
  ΔI Monitor checks for information balance. On convergence, the
  attractor Form is marked as the computation result.

PHASE 6: COMMIT
  results and full DAG journal written to persistent store.
  Returned to caller with cryptographic receipt.
\end{verbatim}

Phases 1--5 are deterministic and pure. Only Phase 6 has external side
effects.

\hypertarget{determinism-guarantees}{%
\subsection{Determinism Guarantees}\label{determinism-guarantees}}

Determinism is not a property of the Bracket Engine alone; it must be
preserved across the entire stack. The guarantees decompose as:

\begin{longtable}[]{@{}
  >{\raggedright\arraybackslash}p{(\columnwidth - 2\tabcolsep) * \real{0.2414}}
  >{\raggedright\arraybackslash}p{(\columnwidth - 2\tabcolsep) * \real{0.7586}}@{}}
\toprule\noalign{}
\begin{minipage}[b]{\linewidth}\raggedright
Layer
\end{minipage} & \begin{minipage}[b]{\linewidth}\raggedright
Determinism Mechanism
\end{minipage} \\
\midrule\noalign{}
\endhead
\bottomrule\noalign{}
\endlastfoot
Ingest & Pydantic's deterministic JSON schema + sorted-keys canonical
form \\
Plan & Symbolic --- no numerical evaluation \\
Execute & SymPy exact rationals; NumPy only with fixed seed and fixed
library versions \\
Validate & Closure defect computed in exact arithmetic \\
Converge & ΔI threshold is a rational constant \\
Commit & Content-addressed store --- same inputs produce same storage
keys \\
\end{longtable}

Floating-point non-determinism is quarantined to the NumPy backend and
is \emph{advisory only}. No theorem-grade output can depend on NumPy.

\hypertarget{reproducibility-envelope}{%
\subsection{Reproducibility Envelope}\label{reproducibility-envelope}}

A computation is reproducible if and only if:

\begin{itemize}
\tightlist
\item
  All inputs are content-addressed (SHA-256)
\item
  All library versions are pinned (\texttt{pyproject.toml} with exact
  hashes)
\item
  The computation grade is TheoremGrade or VerificationGrade
\item
  The computation completes within the BCH-3 envelope (no unbounded
  loops permitted)
\end{itemize}

The reproducibility envelope is explicit: every output carries a
\textbf{receipt} containing input hashes, library hashes, grade, and
timing bounds. Re-execution with the same receipt MUST produce
byte-identical output or the receipt is considered invalidated.

\begin{center}\rule{0.5\linewidth}{0.5pt}\end{center}

\hypertarget{verification-stack}{%
\section{Verification Stack}\label{verification-stack}}

\hypertarget{the-three-tier-verification-model}{%
\subsection{The Three-Tier Verification
Model}\label{the-three-tier-verification-model}}

Verification in the ACS stack is layered to match the BCH truncation:

\textbf{Tier 0: Type-level verification.} Static checks via mypy,
Pydantic schemas, and runtime invariant assertions. Catches schema
violations and codependence breaks before any computation runs.

\textbf{Tier 1: Unit verification.} Each Function's \texttt{apply},
\texttt{bracket\_with}, and \texttt{killing\_form} methods are
unit-tested against golden values. Every algebra element has a test
suite that verifies its bracket relations against the structure
constants.

\textbf{Tier 2: Theorem verification.} The 76-script ACS verification
suite runs as integration tests. Each script corresponds to a specific
theorem (T1--T21) or derived match (D1--D11). A CI run must pass all 76
scripts before a release is cut.

\hypertarget{continuous-verification}{%
\subsection{Continuous Verification}\label{continuous-verification}}

The CI pipeline is layered:

\begin{verbatim}
+-----------------------------------------------------------------+
| TIER 0 (seconds)                                                |
|  mypy · pydantic schema check · invariant assertions            |
+-----------------------------------------------------------------+
| TIER 1 (minutes)                                                |
|  unit tests · bracket relations · Jacobi identity sampling      |
+-----------------------------------------------------------------+
| TIER 2 (hours)                                                  |
|  76-script ACS suite · closure attractor 50k-sample regression |
+-----------------------------------------------------------------+
| RELEASE GATE                                                    |
|  All three tiers green · receipt hash matches previous release |
+-----------------------------------------------------------------+
\end{verbatim}

\hypertarget{two-phase-commit-protocol-for-the-dag-journal}{%
\subsection{Two-Phase Commit Protocol (for the DAG
Journal)}\label{two-phase-commit-protocol-for-the-dag-journal}}

The Bracket Engine must guarantee that no output is observable before
its provenance is journaled. This requires a 2PC protocol:

\begin{enumerate}
\def\labelenumi{\arabic{enumi}.}
\tightlist
\item
  \textbf{Prepare.} Compute the output and stage the journal entry.
\item
  \textbf{Write.} Append the journal entry to the immutable log.
\item
  \textbf{Publish.} Make the output visible to callers.
\item
  \textbf{Rollback path.} If any step fails, the staged journal entry is
  discarded, the output is not published, and the caller receives the
  error atomically.
\end{enumerate}

The journal uses an append-only log structure with Merkle-tree
checkpointing every 1024 entries. Integrity is verifiable offline.

\begin{center}\rule{0.5\linewidth}{0.5pt}\end{center}

\hypertarget{governance-model}{%
\section{Governance Model}\label{governance-model}}

\hypertarget{the-constraint-attractor-cycle}{%
\subsection{The Constraint-Attractor
Cycle}\label{the-constraint-attractor-cycle}}

The governance layer implements the inversion arc from Paper C
(\textit{The Inversion Arc}): every solution, once it satisfies the closure and
information-balance criteria, is promoted to become a constraint for
subsequent computations. This prevents drift and enforces cumulative
consistency.

\begin{verbatim}
   +----------+   solve    +------------+  promote   +----------+
   | CONSTRAINT| --------→ |  ATTRACTOR | ---------→ |CONSTRAINT|
   |   set k   |           |   state    |            |  set k+1 |
   +----------+            +------------+            +----------+
         ^                                                  |
         +---------------- feedback ------------------------+
\end{verbatim}

\hypertarget{grade-hierarchy}{%
\subsection{Grade Hierarchy}\label{grade-hierarchy}}

Outputs are tagged with a grade that determines their usage rights:

\begin{longtable}[]{@{}
  >{\raggedright\arraybackslash}p{(\columnwidth - 6\tabcolsep) * \real{0.1667}}
  >{\raggedright\arraybackslash}p{(\columnwidth - 6\tabcolsep) * \real{0.2143}}
  >{\raggedright\arraybackslash}p{(\columnwidth - 6\tabcolsep) * \real{0.4524}}
  >{\raggedright\arraybackslash}p{(\columnwidth - 6\tabcolsep) * \real{0.1667}}@{}}
\toprule\noalign{}
\begin{minipage}[b]{\linewidth}\raggedright
Grade
\end{minipage} & \begin{minipage}[b]{\linewidth}\raggedright
Backend
\end{minipage} & \begin{minipage}[b]{\linewidth}\raggedright
Closure Threshold
\end{minipage} & \begin{minipage}[b]{\linewidth}\raggedright
Usage
\end{minipage} \\
\midrule\noalign{}
\endhead
\bottomrule\noalign{}
\endlastfoot
Theorem & SymPy (exact ℚ) & \texttt{\textless{}\ 1e-14} & May be cited
as a theorem, added to the permanent constraint set \\
Verification & mpmath & \texttt{\textless{}\ 1e-10} & Strong numerical
evidence; candidate for promotion to Theorem \\
Exploratory & NumPy (fp64) & \texttt{\textless{}\ 1e-6} & Advisory only;
never admitted to the constraint set \\
\end{longtable}

Promotion from Verification to Theorem requires re-running the
computation in SymPy exact arithmetic and passing the stricter
threshold. Downgrade never happens automatically; it requires an
explicit governance action with audit trail.

\hypertarget{audit-log}{%
\subsection{Audit Log}\label{audit-log}}

Every state transition in the governance layer is logged to an
append-only audit store. The log schema:

\begin{verbatim}
{
  "timestamp": "ISO-8601",
  "actor": "string (user or service identity)",
  "action": "ingest | bracket | validate | promote | downgrade | reject",
  "input_hashes": ["sha256:..."],
  "output_hash": "sha256:...",
  "grade": "theorem | verification | exploratory",
  "policy_version": "semver",
  "witness": "string (reproduction instructions)"
}
\end{verbatim}

Logs are cryptographically linked via a hash chain; any tampering breaks
the chain and is detectable in O(N) verification.

\hypertarget{policy-as-code}{%
\subsection{Policy as Code}\label{policy-as-code}}

All governance policies are declarative and version-controlled:

\begin{Shaded}
\begin{Highlighting}[]
\CommentTok{\# policies/theorem\_grade.yaml}
\FunctionTok{name}\KeywordTok{:}\AttributeTok{ TheoremGrade}
\FunctionTok{version}\KeywordTok{:}\AttributeTok{ }\FloatTok{1.2.0}
\FunctionTok{constraints}\KeywordTok{:}
\AttributeTok{  }\KeywordTok{{-}}\AttributeTok{ }\FunctionTok{kind}\KeywordTok{:}\AttributeTok{ backend}
\AttributeTok{    }\FunctionTok{allowed}\KeywordTok{:}\AttributeTok{ }\KeywordTok{[}\AttributeTok{sympy}\KeywordTok{]}
\AttributeTok{  }\KeywordTok{{-}}\AttributeTok{ }\FunctionTok{kind}\KeywordTok{:}\AttributeTok{ closure\_threshold}
\AttributeTok{    }\FunctionTok{max}\KeywordTok{:}\AttributeTok{ }\StringTok{"1e{-}14"}
\AttributeTok{  }\KeywordTok{{-}}\AttributeTok{ }\FunctionTok{kind}\KeywordTok{:}\AttributeTok{ jacobi\_verification}
\AttributeTok{    }\FunctionTok{required}\KeywordTok{:}\AttributeTok{ }\CharTok{true}
\AttributeTok{    }\FunctionTok{sampling}\KeywordTok{:}\AttributeTok{ exhaustive}
\AttributeTok{  }\KeywordTok{{-}}\AttributeTok{ }\FunctionTok{kind}\KeywordTok{:}\AttributeTok{ bch\_order}
\AttributeTok{    }\FunctionTok{max}\KeywordTok{:}\AttributeTok{ }\DecValTok{3}
\FunctionTok{attractor\_criteria}\KeywordTok{:}
\AttributeTok{  }\KeywordTok{{-}}\AttributeTok{ }\FunctionTok{kind}\KeywordTok{:}\AttributeTok{ delta\_I}
\AttributeTok{    }\FunctionTok{threshold}\KeywordTok{:}\AttributeTok{ }\StringTok{"1e{-}12"}
\AttributeTok{    }\FunctionTok{window}\KeywordTok{:}\AttributeTok{ }\DecValTok{8}
\AttributeTok{  }\KeywordTok{{-}}\AttributeTok{ }\FunctionTok{kind}\KeywordTok{:}\AttributeTok{ closure\_stability}
\AttributeTok{    }\FunctionTok{consecutive\_passes}\KeywordTok{:}\AttributeTok{ }\DecValTok{3}
\end{Highlighting}
\end{Shaded}

Policy changes require code review, version bump, and a migration test
--- running the previous theorem set under the new policy to confirm no
regressions.

\begin{center}\rule{0.5\linewidth}{0.5pt}\end{center}

\hypertarget{deployment-architecture}{%
\section{Deployment Architecture}\label{deployment-architecture}}

\hypertarget{deployment-topology}{%
\subsection{Deployment Topology}\label{deployment-topology}}

\begin{verbatim}
+-----------------------------------------------------------------+
|  EDGE — Jupyter kernels · CLI · REST clients                    |
+-----------------------------+-----------------------------------+
                              | TLS 1.3, mTLS for service calls
+-----------------------------v-----------------------------------+
|  API GATEWAY — auth, rate limits, grade routing                 |
+-----------------------------+-----------------------------------+
                              |
+-----------------------------v-----------------------------------+
|  APPLICATION TIER                                               |
|  Bracket Engine workers (stateless) — horizontal scale          |
|  Governance Policy Engine (stateful, replicated 3-way)          |
+--------------+--------------------------+-----------------------+
               |                          |
      +--------v--------+         +------v-------+
      |  Content Store  |         | Audit Log    |
      |  (S3-compatible,|         | (append-only,|
      |  content-addr.) |         | hash-chained)|
      +-----------------+         +--------------+
\end{verbatim}

\hypertarget{service-decomposition}{%
\subsection{Service Decomposition}\label{service-decomposition}}

Four services, clearly bounded:

\begin{enumerate}
\def\labelenumi{\arabic{enumi}.}
\tightlist
\item
  \textbf{\texttt{acs-api}} --- stateless REST façade. Validates inputs,
  routes to the Bracket Engine, returns results with receipts.
\item
  \textbf{\texttt{acs-engine}} --- the Bracket Engine worker pool.
  Consumes work from a deterministic work queue. Stateless modulo the
  journal.
\item
  \textbf{\texttt{acs-governance}} --- Policy Engine, Closure Validator,
  ΔI Monitor. Replicated 3-way with Raft consensus for policy updates.
\item
  \textbf{\texttt{acs-journal}} --- append-only audit + DAG log.
  Content-addressed store. Periodic Merkle checkpoints.
\end{enumerate}

Each service is independently deployable, independently versioned, and
communicates only through explicit message contracts (protobuf).

\hypertarget{resource-envelope-phase-1-target}{%
\subsection{Resource Envelope (Phase 1
Target)}\label{resource-envelope-phase-1-target}}

\begin{longtable}[]{@{}ll@{}}
\toprule\noalign{}
Resource & Sizing \\
\midrule\noalign{}
\endhead
\bottomrule\noalign{}
\endlastfoot
Engine workers & 4 cores, 16 GB each, 3--8 replicas \\
Governance & 2 cores, 8 GB, 3-way replicated \\
API façade & 2 cores, 4 GB, 3-way replicated \\
Journal & 500 GB SSD (+ cold storage) \\
Content store & 5 TB (S3-compatible) \\
\end{longtable}

Phase 1 targets a single region deployment. Multi-region is a Phase 3
concern and requires the content store to support regional
content-addressed replication.

\hypertarget{observability}{%
\subsection{Observability}\label{observability}}

\begin{itemize}
\tightlist
\item
  \textbf{Metrics.} Per-bracket latency, closure-defect distribution, ΔI
  trajectory, grade distribution. Prometheus-format.
\item
  \textbf{Tracing.} Every computation gets a trace ID propagated through
  all services. OpenTelemetry.
\item
  \textbf{Logs.} Structured JSON. Audit log is separate from operational
  log and is never deleted.
\item
  \textbf{Alerts.} Closure threshold violations in Theorem-grade work;
  journal integrity failures; unusual latency distributions in the
  Bracket Engine.
\end{itemize}

\begin{center}\rule{0.5\linewidth}{0.5pt}\end{center}

\hypertarget{api-and-integration-contracts}{%
\section{API and Integration
Contracts}\label{api-and-integration-contracts}}

\hypertarget{rest-api-minimal-surface}{%
\subsection{REST API (Minimal Surface)}\label{rest-api-minimal-surface}}

\begin{verbatim}
POST /v1/bracket
  Body: { form: FormRef, function: FunctionRef, order: 1|2|3, grade: string }
  Returns: { output: FormRef, receipt: string, grade_achieved: string }

GET /v1/form/{content_hash}
  Returns the full Form by its content hash.

POST /v1/validate
  Body: { form: FormRef, subspace: AlgebraSubspaceRef, threshold: Decimal }
  Returns: { passed: bool, defect: Decimal, witness: string }

POST /v1/verify/{theorem_id}
  Runs the verification script for the named theorem.
  Returns: { passed: bool, execution_time: duration, receipt: string }

GET /v1/audit/{receipt}
  Returns the complete audit trail for a given receipt.
\end{verbatim}

All endpoints are idempotent on the content-hash identity. No session
state. No conversational memory.

\hypertarget{python-client-api}{%
\subsection{Python Client API}\label{python-client-api}}

\begin{Shaded}
\begin{Highlighting}[]
\ImportTok{from}\NormalTok{ acs }\ImportTok{import}\NormalTok{ Stack, Form, Function}

\NormalTok{stack }\OperatorTok{=}\NormalTok{ Stack.}\ExtensionTok{connect}\NormalTok{(}\StringTok{"https://acs.example.org"}\NormalTok{)}

\CommentTok{\# Load a Form}
\NormalTok{f }\OperatorTok{=}\NormalTok{ stack.forms.get(}\StringTok{"sha256:abc123..."}\NormalTok{)}

\CommentTok{\# Apply a Function}
\NormalTok{phi }\OperatorTok{=}\NormalTok{ stack.functions.get(}\StringTok{"bracket.T\_BL.sl4"}\NormalTok{)}

\CommentTok{\# Compose — uses the | operator for bracket}
\NormalTok{result }\OperatorTok{=}\NormalTok{ f }\OperatorTok{|}\NormalTok{ phi  }\CommentTok{\# implicit order 3, Verification grade}

\CommentTok{\# Explicit version with full control}
\NormalTok{result }\OperatorTok{=}\NormalTok{ stack.bracket(}
\NormalTok{    form}\OperatorTok{=}\NormalTok{f,}
\NormalTok{    function}\OperatorTok{=}\NormalTok{phi,}
\NormalTok{    order}\OperatorTok{=}\DecValTok{3}\NormalTok{,}
\NormalTok{    grade}\OperatorTok{=}\StringTok{"Theorem"}\NormalTok{,}
\NormalTok{)}

\CommentTok{\# Every result carries its receipt}
\BuiltInTok{print}\NormalTok{(result.receipt)}
\CommentTok{\# {-}\textgreater{} "receipt:sha256:..., grade:Theorem, policy:v1.2.0, ts:2026{-}04{-}17T..."}
\end{Highlighting}
\end{Shaded}

\hypertarget{integration-with-external-tools}{%
\subsection{Integration with External
Tools}\label{integration-with-external-tools}}

The stack ships adapters for:

\begin{itemize}
\tightlist
\item
  \textbf{Jupyter} --- a kernel that exposes the REST API through
  \texttt{\%acs} magic
\item
  \textbf{SymPy} --- round-trip conversion between SymPy expressions and
  ACS Forms
\item
  \textbf{LaTeX} --- theorem output emits LaTeX with cryptographic
  provenance footnotes linking back to the receipt
\item
  \textbf{Git} --- the audit log can be mirrored into a git-compatible
  append-only branch for version-controlled collaboration
\end{itemize}

\begin{center}\rule{0.5\linewidth}{0.5pt}\end{center}

\hypertarget{phased-implementation-roadmap}{%
\section{Phased Implementation
Roadmap}\label{phased-implementation-roadmap}}

\hypertarget{phase-0-foundations-months-02}{%
\subsection{Phase 0 --- Foundations (Months
0--2)}\label{phase-0-foundations-months-02}}

\begin{itemize}
\tightlist
\item
  Kernel (Layer 0): SymPy + mpmath + NumPy backends, kernel invariant
  test suite
\item
  Form/Function (Layer 1): Pydantic schemas, type system, 20+ standard
  Form types and Function types
\item
  Content-addressed store prototype
\item
  Audit log prototype
\end{itemize}

\textbf{Exit criterion.} All 76 verification scripts pass under the new
Form/Function typing.

\hypertarget{phase-1-core-engine-months-25}{%
\subsection{Phase 1 --- Core Engine (Months
2--5)}\label{phase-1-core-engine-months-25}}

\begin{itemize}
\tightlist
\item
  Bracket Engine with full BCH-3 support
\item
  Closure Validator
\item
  ΔI Monitor
\item
  Basic Policy Engine (TheoremGrade, VerificationGrade,
  ExploratoryGrade)
\item
  REST API façade
\end{itemize}

\textbf{Exit criterion.} End-to-end determinism test: given a fixed Form
and Function, the stack produces byte-identical outputs across 1000
consecutive invocations, and across clean-room re-deployments.

\hypertarget{phase-2-governance-months-58}{%
\subsection{Phase 2 --- Governance (Months
5--8)}\label{phase-2-governance-months-58}}

\begin{itemize}
\tightlist
\item
  Full constraint-attractor cycle
\item
  Policy-as-code with version migration tests
\item
  Audit log Merkle checkpointing
\item
  Grade promotion workflow with human-in-the-loop approval
\end{itemize}

\textbf{Exit criterion.} Independent auditor can replay any
Theorem-grade output from its receipt alone and confirm byte-identity.

\hypertarget{phase-3-scale-and-integration-months-812}{%
\subsection{Phase 3 --- Scale and Integration (Months
8--12)}\label{phase-3-scale-and-integration-months-812}}

\begin{itemize}
\tightlist
\item
  Horizontal scaling of the Bracket Engine
\item
  Multi-region journal replication
\item
  Jupyter + SymPy + LaTeX adapters
\item
  Public API with authenticated access
\end{itemize}

\textbf{Exit criterion.} 99.9\% uptime under production load; sub-second
latency for bracket operations below complexity bound C₁ (tbd).

\hypertarget{phase-4-extension-to-the-physics-frontier-months-12}{%
\subsection{Phase 4 --- Extension to the Physics Frontier (Months
12+)}\label{phase-4-extension-to-the-physics-frontier-months-12}}

\begin{itemize}
\tightlist
\item
  Lattice PS simulation integration (external compute cluster)
\item
  Neutrino precision posterior fitter
\item
  GW waveform templates using the 0:1:4 torsion hierarchy
\item
  X-ray line search posterior updater (XRISM data feed)
\end{itemize}

\textbf{Exit criterion.} Phase 4 is open-ended; each experimental test
from Paper A's Appendix E maps to a service integration milestone.

\begin{center}\rule{0.5\linewidth}{0.5pt}\end{center}

\hypertarget{risks-and-mitigations}{%
\section{Risks and Mitigations}\label{risks-and-mitigations}}

\begin{longtable}[]{@{}
  >{\raggedright\arraybackslash}p{(\columnwidth - 6\tabcolsep) * \real{0.1500}}
  >{\raggedright\arraybackslash}p{(\columnwidth - 6\tabcolsep) * \real{0.3000}}
  >{\raggedright\arraybackslash}p{(\columnwidth - 6\tabcolsep) * \real{0.2500}}
  >{\raggedright\arraybackslash}p{(\columnwidth - 6\tabcolsep) * \real{0.3000}}@{}}
\toprule\noalign{}
\begin{minipage}[b]{\linewidth}\raggedright
Risk
\end{minipage} & \begin{minipage}[b]{\linewidth}\raggedright
Likelihood
\end{minipage} & \begin{minipage}[b]{\linewidth}\raggedright
Severity
\end{minipage} & \begin{minipage}[b]{\linewidth}\raggedright
Mitigation
\end{minipage} \\
\midrule\noalign{}
\endhead
\bottomrule\noalign{}
\endlastfoot
Floating-point creep into Theorem-grade work & Medium & Critical &
Policy Engine refuses to grade as Theorem unless SymPy backend is used;
CI enforces \\
Journal corruption & Low & Critical & Merkle checkpoints every 1024
entries; offline integrity verifier; 3-way replication \\
Policy version drift between services & Medium & High & All services
embed the policy version in every response; mismatch aborts the call \\
Performance bottleneck at BCH-3 & High & Medium & BCH-3 is algebraically
bounded; if dominant, introduce lazy evaluation and memoisation.
Complexity is O(k³) in algebra rank, not exponential \\
Closure threshold tuning & Medium & Medium & Thresholds are part of
policy-as-code; changes require migration tests against the full theorem
set \\
Jacobi identity violations from numerical error & Medium & High & Tier-1
CI samples 10⁴ random triples per build; violation triggers immediate
policy-engine quarantine \\
Determinism break from library upgrade & High & Medium & All
dependencies pinned by hash; upgrade requires full 76-script
regression \\
Abuse of Exploratory grade for production claims & Low & High & Grade is
cryptographically bound to receipt; external tools reject unapproved
grades \\
Auditor-replay failure on legitimate output & Low & Critical & Receipts
carry complete reproduction envelope including library hashes; replay
tools ship with the receipt format spec \\
\end{longtable}

\begin{center}\rule{0.5\linewidth}{0.5pt}\end{center}

\hypertarget{traceability-to-the-acs-trilogy}{%
\section{Traceability to the ACS
Trilogy}\label{traceability-to-the-acs-trilogy}}

Every architectural decision in this document traces back to a specific
result in the ACS trilogy. This section is the cross-reference index.

\begin{longtable}[]{@{}
  >{\raggedright\arraybackslash}p{(\columnwidth - 4\tabcolsep) * \real{0.4250}}
  >{\raggedright\arraybackslash}p{(\columnwidth - 4\tabcolsep) * \real{0.3750}}
  >{\raggedright\arraybackslash}p{(\columnwidth - 4\tabcolsep) * \real{0.2000}}@{}}
\toprule\noalign{}
\begin{minipage}[b]{\linewidth}\raggedright
Design Decision
\end{minipage} & \begin{minipage}[b]{\linewidth}\raggedright
Justification
\end{minipage} & \begin{minipage}[b]{\linewidth}\raggedright
Source
\end{minipage} \\
\midrule\noalign{}
\endhead
\bottomrule\noalign{}
\endlastfoot
Form/Function type split & ACS-2 structural asymmetry & Paper A §2.1 Def
2.3 \\
BCH-3 truncation & Jacobi truncation theorem & Paper A §5.1 Thm \\
Closure-defect validation & Closure attractor selection (50k samples) &
Paper A §4.3 \\
Information-balance halting & \texttt{ΔI\ =\ 0} attractor & Paper A §2.3
Lemma 2.10 \\
Constraint-attractor cycle & Inversion arc & Paper C §3 \\
Content-addressed provenance & Reproducibility envelope (Paper A
footnote) & Paper A §6.8 \\
Three grades (Theorem/Ver/Exp) & Epistemic tiers & Paper A §C.2 \\
Torsion-hierarchy priority weights & 0:1:4 coupling hierarchy & Paper A
§C.1 \\
Seven configuration parameters & 7-parameter irreducible boundary &
Paper A §C.2 \\
Two-phase commit & Constraint before publication & Paper C §3.3 \\
\end{longtable}

\begin{center}\rule{0.5\linewidth}{0.5pt}\end{center}

\hypertarget{glossary}{%
\section{Glossary}\label{glossary}}

\textbf{ACS.} Asymmetric Codependent System --- a pair of fields (Form,
Function) that cannot be defined independently and whose interaction is
governed by a Lie bracket.

\textbf{BCH.} Baker-Campbell-Hausdorff formula --- the expansion for
\texttt{log(exp(X)\ exp(Y))} as a series of nested commutators.

\textbf{Bracket Engine.} The core compositional primitive of the stack.
Computes \texttt{{[}F,\ Φ{]}} truncated at BCH order 3.

\textbf{Closure attractor.} The unique subspace of a Lie algebra in
which a given set of generators closes under the bracket, minimising the
closure defect functional.

\textbf{Closure defect \texttt{D(V)}.} A functional measuring how far a
subspace \texttt{V} falls short of being closed under the bracket.

\textbf{Constraint-attractor cycle.} The governance feedback loop in
which outputs that satisfy the current constraints become the
constraints for subsequent computations.

\textbf{Form.} A typed immutable state object. Corresponds to the
vierbein in the physics.

\textbf{Function.} A typed operator acting on Forms. Corresponds to the
spin connection in the physics.

\textbf{Grade.} A classification of computational trust level: Theorem,
Verification, or Exploratory.

\textbf{Holonomy.} The third-order BCH term
\texttt{{[}{[}f,g{]},\ h{]}}. In the architecture, this is the final
emergent layer before Jacobi truncation.

\textbf{Information balance.} The condition \texttt{ΔI\ =\ 0} indicating
a computation has reached its attractor.

\textbf{Jacobi truncation.} The fact that
\texttt{{[}{[}A,B{]},C{]}\ +\ {[}{[}B,C{]},A{]}\ +\ {[}{[}C,A{]},B{]}\ =\ 0}
ensures no new independent content is generated beyond BCH order 3.

\textbf{Killing form.} The symmetric bilinear form
\texttt{K(X,\ Y)\ =\ tr(ad(X)\ ad(Y))} on a Lie algebra.

\textbf{Receipt.} A cryptographic record of a computation, sufficient
for independent replay.

\begin{center}\rule{0.5\linewidth}{0.5pt}\end{center}

\hypertarget{appendix-a-reference-implementation-sketch}{%
\section{Appendix A --- Reference Implementation
Sketch}\label{appendix-a-reference-implementation-sketch}}

The following skeleton shows the core of the Bracket Engine in minimal
Python form, demonstrating determinism and provenance.

\begin{Shaded}
\begin{Highlighting}[]
\CommentTok{\# acs/engine.py}
\ImportTok{from}\NormalTok{ \_\_future\_\_ }\ImportTok{import}\NormalTok{ annotations}
\ImportTok{from}\NormalTok{ dataclasses }\ImportTok{import}\NormalTok{ dataclass, replace}
\ImportTok{from}\NormalTok{ hashlib }\ImportTok{import}\NormalTok{ sha256}
\ImportTok{from}\NormalTok{ decimal }\ImportTok{import}\NormalTok{ Decimal}
\ImportTok{from}\NormalTok{ typing }\ImportTok{import}\NormalTok{ Protocol}
\ImportTok{import}\NormalTok{ json}

\KeywordTok{class}\NormalTok{ FormProto(Protocol):}
\NormalTok{    content: }\BuiltInTok{object}
\NormalTok{    schema\_version: }\BuiltInTok{str}
\NormalTok{    provenance: }\BuiltInTok{str}
    \AttributeTok{@property}
    \KeywordTok{def}\NormalTok{ content\_hash(}\VariableTok{self}\NormalTok{) }\OperatorTok{{-}\textgreater{}} \BuiltInTok{str}\NormalTok{: ...}

\KeywordTok{class}\NormalTok{ FunctionProto(Protocol):}
\NormalTok{    name: }\BuiltInTok{str}
\NormalTok{    order: }\BuiltInTok{int}
    \KeywordTok{def} \BuiltInTok{apply}\NormalTok{(}\VariableTok{self}\NormalTok{, f: FormProto) }\OperatorTok{{-}\textgreater{}}\NormalTok{ FormProto: ...}
    \KeywordTok{def}\NormalTok{ bracket\_with(}\VariableTok{self}\NormalTok{, other: FunctionProto) }\OperatorTok{{-}\textgreater{}}\NormalTok{ FunctionProto: ...}

\AttributeTok{@dataclass}\NormalTok{(frozen}\OperatorTok{=}\VariableTok{True}\NormalTok{)}
\KeywordTok{class}\NormalTok{ Receipt:}
\NormalTok{    input\_hash: }\BuiltInTok{str}
\NormalTok{    function\_name: }\BuiltInTok{str}
\NormalTok{    order: }\BuiltInTok{int}
\NormalTok{    output\_hash: }\BuiltInTok{str}
\NormalTok{    policy\_version: }\BuiltInTok{str}
\NormalTok{    grade: }\BuiltInTok{str}

\NormalTok{\_JOURNAL: }\BuiltInTok{list}\NormalTok{[Receipt] }\OperatorTok{=}\NormalTok{ []  }\CommentTok{\# Prototype — real impl uses append{-}only store}

\KeywordTok{def}\NormalTok{ bracket(}
\NormalTok{    f: FormProto,}
\NormalTok{    phi: FunctionProto,}
\NormalTok{    order: }\BuiltInTok{int} \OperatorTok{=} \DecValTok{3}\NormalTok{,}
\NormalTok{    grade: }\BuiltInTok{str} \OperatorTok{=} \StringTok{"Verification"}\NormalTok{,}
\NormalTok{    policy\_version: }\BuiltInTok{str} \OperatorTok{=} \StringTok{"1.0.0"}\NormalTok{,}
\NormalTok{) }\OperatorTok{{-}\textgreater{}}\NormalTok{ FormProto:}
    \ControlFlowTok{if}\NormalTok{ order }\KeywordTok{not} \KeywordTok{in}\NormalTok{ (}\DecValTok{1}\NormalTok{, }\DecValTok{2}\NormalTok{, }\DecValTok{3}\NormalTok{):}
        \ControlFlowTok{raise} \PreprocessorTok{ValueError}\NormalTok{(}\SpecialStringTok{f"BCH order must be 1, 2, or 3; got }\SpecialCharTok{\{}\NormalTok{order}\SpecialCharTok{\}}\SpecialStringTok{"}\NormalTok{)}

    \CommentTok{\# Phase: execute the truncated bracket}
\NormalTok{    r1 }\OperatorTok{=}\NormalTok{ phi.}\BuiltInTok{apply}\NormalTok{(f)}
    \ControlFlowTok{if}\NormalTok{ order }\OperatorTok{==} \DecValTok{1}\NormalTok{:}
\NormalTok{        out }\OperatorTok{=}\NormalTok{ r1}
    \ControlFlowTok{else}\NormalTok{:}
\NormalTok{        r2 }\OperatorTok{=}\NormalTok{ \_commutator(f, r1)}
        \ControlFlowTok{if}\NormalTok{ order }\OperatorTok{==} \DecValTok{2}\NormalTok{:}
\NormalTok{            out }\OperatorTok{=}\NormalTok{ r2}
        \ControlFlowTok{else}\NormalTok{:  }\CommentTok{\# order == 3}
\NormalTok{            r3a }\OperatorTok{=}\NormalTok{ \_commutator(f, r2)}
\NormalTok{            r3b }\OperatorTok{=}\NormalTok{ \_commutator(r1, r2)}
\NormalTok{            out }\OperatorTok{=}\NormalTok{ \_add(r3a, r3b)}

    \CommentTok{\# Phase: validate closure (elided here)}
    \CommentTok{\# Phase: compute provenance}
\NormalTok{    new\_provenance }\OperatorTok{=}\NormalTok{ sha256(}
        \SpecialStringTok{f"}\SpecialCharTok{\{}\NormalTok{f}\SpecialCharTok{.}\NormalTok{content\_hash}\SpecialCharTok{\}}\SpecialStringTok{|}\SpecialCharTok{\{}\NormalTok{phi}\SpecialCharTok{.}\NormalTok{name}\SpecialCharTok{\}}\SpecialStringTok{|}\SpecialCharTok{\{}\NormalTok{order}\SpecialCharTok{\}}\SpecialStringTok{"}\NormalTok{.encode()}
\NormalTok{    ).hexdigest()}
\NormalTok{    out }\OperatorTok{=}\NormalTok{ replace(out, provenance}\OperatorTok{=}\SpecialStringTok{f"sha256:}\SpecialCharTok{\{}\NormalTok{new\_provenance}\SpecialCharTok{\}}\SpecialStringTok{"}\NormalTok{)}

    \CommentTok{\# Phase: journal}
\NormalTok{    \_JOURNAL.append(Receipt(}
\NormalTok{        input\_hash}\OperatorTok{=}\NormalTok{f.content\_hash,}
\NormalTok{        function\_name}\OperatorTok{=}\NormalTok{phi.name,}
\NormalTok{        order}\OperatorTok{=}\NormalTok{order,}
\NormalTok{        output\_hash}\OperatorTok{=}\NormalTok{out.content\_hash,}
\NormalTok{        policy\_version}\OperatorTok{=}\NormalTok{policy\_version,}
\NormalTok{        grade}\OperatorTok{=}\NormalTok{grade,}
\NormalTok{    ))}
    \ControlFlowTok{return}\NormalTok{ out}

\KeywordTok{def}\NormalTok{ \_commutator(a: FormProto, b: FormProto) }\OperatorTok{{-}\textgreater{}}\NormalTok{ FormProto:}
    \CommentTok{\# Implemented per algebra type; the kernel dispatches on schema\_version.}
\NormalTok{    ...}

\KeywordTok{def}\NormalTok{ \_add(a: FormProto, b: FormProto) }\OperatorTok{{-}\textgreater{}}\NormalTok{ FormProto:}
\NormalTok{    ...}
\end{Highlighting}
\end{Shaded}

This skeleton satisfies invariants I-B1 (same inputs → same outputs via
deterministic \texttt{phi.apply}), I-B2 (BCH capped at 3), and I-B3
(journal append happens before return). A production implementation
replaces the in-memory \texttt{\_JOURNAL} with a proper append-only
store and adds the 2PC protocol from §5.3.

\begin{center}\rule{0.5\linewidth}{0.5pt}\end{center}

\hypertarget{document-status}{%
\section{Document Status}\label{document-status}}

\begin{itemize}
\tightlist
\item
  \textbf{Version.} 1.0 (Initial PDR)
\item
  \textbf{Status.} Design review --- open for technical critique
\item
  \textbf{Next review.} After Phase 0 exit criterion
\item
  \textbf{Author.} Bradley Wallace
\item
  \textbf{Document ID.} ACS Deterministic AI Stack PDR
\end{itemize}

\begin{center}\rule{0.5\linewidth}{0.5pt}\end{center}

\emph{End of document.}

\end{document}
